\section{解谱法}

\subsection{有限维线性解谱问题}

设真实直方图为
\begin{equation*}
	\vb{\mu}=(\mu_1,\mu_2,\ldots,\mu_M)^T,
\end{equation*}
观测直方图为
\begin{equation*}
	\vb{n}=(n_1,n_2,\ldots,n_N)^T.
\end{equation*}
响应矩阵 \(R\in\mathbb{R}^{N\times M}\) 定义为
\begin{equation*}
	\nu_i=\sum_{j=1}^{M}R_{ij}\mu_j+b_i,
	\quad i=1,2,\ldots,N,
\end{equation*}
也即
\begin{equation*}
	\vb{\nu}=R\vb{\mu}+\vb{b}.
\end{equation*}
若本底 \(\vb{b}\) 已知，记
\begin{equation*}
	\vb{y}=\vb{n}-\vb{b}.
\end{equation*}
以下只讨论能写出闭式解的二次型情形.设
\begin{equation*}
	E(\vb{y})=R\vb{\mu}_0,\quad
	\operatorname{cov}(\vb{y})=V,
\end{equation*}
其中 \(V\) 对称正定.记
\begin{equation*}
	W=V^{-1}.
\end{equation*}
加权二次型为
\begin{equation*}
	\chi^2(\vb{\mu})
	=(\vb{y}-R\vb{\mu})^TW(\vb{y}-R\vb{\mu}).
\end{equation*}
考虑最一般的二次正则化闭式问题
\begin{equation*}
	Q(\vb{\mu})
	=\frac{\alpha}{2}(\vb{y}-R\vb{\mu})^TW(\vb{y}-R\vb{\mu})
	+\vb{\mu}^TG\vb{\mu},
\end{equation*}
其中 \(\alpha>0\)，且 \(G=G^T\geq0\).若 \(Q\) 需要在若干线性约束下取极小值，写作
\begin{equation*}
	C^T\vb{\mu}=\vb{d},
\end{equation*}
其中 \(C\in\mathbb{R}^{M\times p}\)，且 \(C\) 的 \(p\) 列线性无关.

\subsection{无约束情形的闭式解}

先不考虑约束.将 \(Q\) 展开为
\begin{equation*}
	Q(\vb{\mu})
	=\frac{\alpha}{2}\vb{y}^TW\vb{y}
	-\alpha\vb{\mu}^TR^TW\vb{y}
	+\frac{1}{2}\vb{\mu}^T\alpha R^TWR\vb{\mu}
	+\vb{\mu}^TG\vb{\mu}.
\end{equation*}
对 \(\vb{\mu}\) 作变分，有
\begin{equation*}
	\delta Q
	=-\alpha\delta\vb{\mu}^TR^TW\vb{y}
	+\alpha\delta\vb{\mu}^TR^TWR\vb{\mu}
	+2\delta\vb{\mu}^TG\vb{\mu}.
\end{equation*}
因此驻点方程为
\begin{equation*}
	\left(\alpha R^TWR+2G\right)\hat{\vb{\mu}}
	=\alpha R^TW\vb{y}.
\end{equation*}
记
\begin{equation*}
	M=\alpha R^TWR+2G.
\end{equation*}
若 \(M\) 非奇异，则无约束闭式解为
\begin{equation*}
	\hat{\vb{\mu}}
	=M^{-1}\alpha R^TW\vb{y}.
\end{equation*}
由于
\begin{equation*}
	\delta^2Q
	=\delta\vb{\mu}^TM\delta\vb{\mu},
\end{equation*}
若 \(M\) 正定，则该驻点是唯一极小值.

\begin{notes}
	若 \(G=0\)，则唯一性要求 \(R^TWR\) 正定，也即 \(R\) 的列满秩.若 \(R\) 不列满秩，则数据中不存在的信息不能由无正则化二次型唯一确定.
\end{notes}

令
\begin{equation*}
	A=M^{-1}\alpha R^TW,
\end{equation*}
则
\begin{equation*}
	\hat{\vb{\mu}}=A\vb{y}.
\end{equation*}
由 \(E(\vb{y})=R\vb{\mu}_0\)，有
\begin{equation*}
	E(\hat{\vb{\mu}})
	=AR\vb{\mu}_0.
\end{equation*}
故偏差为
\begin{equation*}
	\vb{b}_{\rm bias}
	=E(\hat{\vb{\mu}})-\vb{\mu}_0
	=(AR-I)\vb{\mu}_0.
\end{equation*}
利用
\begin{equation*}
	AR-I
	=M^{-1}\alpha R^TWR-I
	=M^{-1}\left(\alpha R^TWR-M\right)
	=-2M^{-1}G,
\end{equation*}
得到
\begin{equation*}
	\vb{b}_{\rm bias}
	=-2M^{-1}G\vb{\mu}_0.
\end{equation*}
协方差为
\begin{equation*}
	U
	=\operatorname{cov}(\hat{\vb{\mu}})
	=A\operatorname{cov}(\vb{y})A^T
	=AVA^T.
\end{equation*}
由于 \(W=V^{-1}\) 且 \(M=M^T\)，可化简为
\begin{equation*}
	U
	=\alpha^2M^{-1}R^TWVWRM^{-1}
	=\alpha^2M^{-1}R^TWRM^{-1}.
\end{equation*}

\subsection{线性约束下的闭式解}

现在考虑约束
\begin{equation*}
	C^T\vb{\mu}=\vb{d}.
\end{equation*}
引入 Lagrange 乘子 \(\vb{\lambda}\in\mathbb{R}^p\)，令
\begin{equation*}
	\mathcal{L}(\vb{\mu},\vb{\lambda})
	=Q(\vb{\mu})+\vb{\lambda}^T(C^T\vb{\mu}-\vb{d}).
\end{equation*}
驻点方程为
\begin{equation*}
	M\hat{\vb{\mu}}+C\hat{\vb{\lambda}}
	=\alpha R^TW\vb{y},
	\quad
	C^T\hat{\vb{\mu}}=\vb{d}.
\end{equation*}
这等价于分块线性方程
\begin{equation*}
	\begin{pmatrix}
		M & C\\
		C^T & 0
	\end{pmatrix}
	\begin{pmatrix}
		\hat{\vb{\mu}}\\
		\hat{\vb{\lambda}}
	\end{pmatrix}
	=
	\begin{pmatrix}
		\alpha R^TW\vb{y}\\
		\vb{d}
	\end{pmatrix}.
\end{equation*}
先记无约束解
\begin{equation*}
	\vb{\mu}_*=
	M^{-1}\alpha R^TW\vb{y}.
\end{equation*}
由第一行得
\begin{equation*}
	\hat{\vb{\mu}}
	=\vb{\mu}_*-M^{-1}C\hat{\vb{\lambda}}.
\end{equation*}
代入约束条件，得到
\begin{equation*}
	C^T\vb{\mu}_*-C^TM^{-1}C\hat{\vb{\lambda}}=\vb{d}.
\end{equation*}
若
\begin{equation*}
	H_C=C^TM^{-1}C
\end{equation*}
非奇异，则
\begin{equation*}
	\hat{\vb{\lambda}}
	=H_C^{-1}(C^T\vb{\mu}_*-\vb{d}).
\end{equation*}
于是线性约束下的闭式解为
\begin{equation*}
	\hat{\vb{\mu}}
	=\vb{\mu}_*
	+M^{-1}C(C^TM^{-1}C)^{-1}
	(\vb{d}-C^T\vb{\mu}_*).
\end{equation*}
这说明约束解等于无约束解加上一个修正项，该修正项属于 \(M^{-1}C\) 张成的子空间，并且恰好使 \(C^T\hat{\vb{\mu}}=\vb{d}\).

\begin{notes}
	若 \(M\) 正定且 \(C\) 列满秩，则 \(C^TM^{-1}C\) 正定.证明为：对任意非零 \(\vb{a}\in\mathbb{R}^p\)，有 \(C\vb{a}\neq0\)，因此
	\begin{equation*}
		\vb{a}^TC^TM^{-1}C\vb{a}
		=(C\vb{a})^TM^{-1}(C\vb{a})>0.
	\end{equation*}
\end{notes}

若约束右端也是数据的线性函数，即
\begin{equation*}
	\vb{d}=D\vb{y},
\end{equation*}
其中 \(D\in\mathbb{R}^{p\times N}\)，则估计量仍是 \(\vb{y}\) 的线性函数.令
\begin{equation*}
	A=M^{-1}\alpha R^TW,
\end{equation*}
则
\begin{equation*}
	\hat{\vb{\mu}}
	=
	A\vb{y}
	+M^{-1}C(C^TM^{-1}C)^{-1}(D-C^TA)\vb{y}.
\end{equation*}
记
\begin{equation*}
	B
	=A
	+M^{-1}C(C^TM^{-1}C)^{-1}(D-C^TA),
\end{equation*}
则
\begin{equation*}
	\hat{\vb{\mu}}=B\vb{y},\quad
	U=BVB^T,\quad
	\vb{b}_{\rm bias}=(BR-I)\vb{\mu}_0.
\end{equation*}

\subsection{总数约束}

常见的总数约束为
\begin{equation*}
	\sum_{i=1}^{N}\nu_i=\sum_{i=1}^{N}n_i.
\end{equation*}
在本底已扣除的记号下，即
\begin{equation*}
	\vb{1}_N^TR\vb{\mu}=\vb{1}_N^T\vb{y}.
\end{equation*}
令
\begin{equation*}
	\vb{c}=R^T\vb{1}_N.
\end{equation*}
则总数约束为
\begin{equation*}
	\vb{c}^T\vb{\mu}=\vb{1}_N^T\vb{y}.
\end{equation*}
这是上一节公式的特例：
\begin{equation*}
	C=\vb{c},\quad D=\vb{1}_N^T.
\end{equation*}
因此
\begin{equation*}
	\hat{\vb{\mu}}
	=
	\vb{\mu}_*
	+M^{-1}\vb{c}\,
	\frac{\vb{1}_N^T\vb{y}-\vb{c}^T\vb{\mu}_*}
	{\vb{c}^TM^{-1}\vb{c}},
\end{equation*}
其中
\begin{equation*}
	\vb{\mu}_*=M^{-1}\alpha R^TW\vb{y}.
\end{equation*}
把它写成 \(\hat{\vb{\mu}}=B\vb{y}\)，则
\begin{equation*}
	B
	=M^{-1}\alpha R^TW
	+M^{-1}\vb{c}\,
	\frac{\vb{1}_N^T-\vb{c}^TM^{-1}\alpha R^TW}
	{\vb{c}^TM^{-1}\vb{c}}.
\end{equation*}
因此
\begin{equation*}
	U_{\rm constr}=BVB^T.
\end{equation*}

\subsection{白化变换与无正则化最小二乘}

现在详细推导线性代数结构.由于 \(V\) 对称正定，可以取唯一的对称正定平方根
\begin{equation*}
	V^{1/2}V^{1/2}=V.
\end{equation*}
记
\begin{equation*}
	\tilde{R}=V^{-1/2}R,\quad
	\tilde{\vb{y}}=V^{-1/2}\vb{y}.
\end{equation*}
则
\begin{equation*}
	(\vb{y}-R\vb{\mu})^TW(\vb{y}-R\vb{\mu})
	=\left\|\tilde{\vb{y}}-\tilde{R}\vb{\mu}\right\|^2.
\end{equation*}
因此无正则化加权最小二乘问题等价于
\begin{equation*}
	\min_{\vb{\mu}}\left\|\tilde{\vb{y}}-\tilde{R}\vb{\mu}\right\|^2.
\end{equation*}
对 \(\tilde{R}\) 作奇异值分解.设
\begin{equation*}
	r=\operatorname{rank}\tilde{R},
\end{equation*}
则存在正交矩阵
\begin{equation*}
	U=(\vb{u}_1,\ldots,\vb{u}_N),\quad
	Q=(\vb{q}_1,\ldots,\vb{q}_M),
\end{equation*}
使得
\begin{equation*}
	\tilde{R}\vb{q}_a=s_a\vb{u}_a,\quad
	\tilde{R}^T\vb{u}_a=s_a\vb{q}_a,\quad
	a=1,\ldots,r,
\end{equation*}
其中
\begin{equation*}
	s_1\geq s_2\geq\cdots\geq s_r>0.
\end{equation*}
把未知量和数据投影到奇异向量基底中：
\begin{equation*}
	\vb{z}=Q^T\vb{\mu},\quad
	\vb{d}=U^T\tilde{\vb{y}}.
\end{equation*}
由于 \(U,Q\) 正交，范数不变，于是
\begin{equation*}
	\left\|\tilde{\vb{y}}-\tilde{R}\vb{\mu}\right\|^2
	=\left\|U^T\tilde{\vb{y}}-U^T\tilde{R}Q\vb{z}\right\|^2.
\end{equation*}
而
\begin{equation*}
	U^T\tilde{R}Q
	=
	\begin{pmatrix}
		s_1 & 0 & \cdots & 0\\
		0 & s_2 & \cdots & 0\\
		\vdots & \vdots & \ddots & \vdots\\
		0 & 0 & \cdots & s_r\\
		0 & 0 & \cdots & 0
	\end{pmatrix}.
\end{equation*}
因此目标函数化为
\begin{equation*}
	\left\|\tilde{\vb{y}}-\tilde{R}\vb{\mu}\right\|^2
	=\sum_{a=1}^{r}(d_a-s_az_a)^2
	+\sum_{a=r+1}^{N}d_a^2,
\end{equation*}
其中当 \(a>r\) 时，若 \(a>M\) 则没有对应的 \(z_a\).

对 \(a=1,\ldots,r\)，极小条件给出
\begin{equation*}
	z_a=\frac{d_a}{s_a}.
\end{equation*}
若 \(r=M\)，则解唯一，且
\begin{equation*}
	\hat{\vb{\mu}}
	=\sum_{a=1}^{M}\frac{d_a}{s_a}\vb{q}_a.
\end{equation*}
若 \(r<M\)，则 \(z_{r+1},\ldots,z_M\) 不被目标函数决定.最小范数解取
\begin{equation*}
	z_{r+1}=\cdots=z_M=0,
\end{equation*}
从而
\begin{equation*}
	\hat{\vb{\mu}}_{\rm min}
	=\sum_{a=1}^{r}\frac{d_a}{s_a}\vb{q}_a.
\end{equation*}
这就是 Moore--Penrose 伪逆给出的解：
\begin{equation*}
	\hat{\vb{\mu}}_{\rm min}
	=\tilde{R}^{+}\tilde{\vb{y}}.
\end{equation*}

\subsection{SVD 基底中的方差放大}

设真实参数在右奇异向量基底中的坐标为
\begin{equation*}
	\vb{\mu}_0=\sum_{a=1}^{M}z_{0a}\vb{q}_a.
\end{equation*}
又令
\begin{equation*}
	\tilde{\vb{y}}=\tilde{R}\vb{\mu}_0+\vb{\epsilon},
	\quad
	E(\vb{\epsilon})=0,\quad
	\operatorname{cov}(\vb{\epsilon})=I_N.
\end{equation*}
投影到 \(\vb{u}_a\) 上，得到
\begin{equation*}
	d_a=\vb{u}_a^T\tilde{\vb{y}}
	=s_az_{0a}+\epsilon_a,
	\quad
	\epsilon_a=\vb{u}_a^T\vb{\epsilon}.
\end{equation*}
由于 \(U\) 正交，
\begin{equation*}
	\operatorname{cov}(\epsilon_a,\epsilon_b)=\delta_{ab}.
\end{equation*}
无正则化反演中
\begin{equation*}
	\hat{z}_a=\frac{d_a}{s_a}
	=z_{0a}+\frac{\epsilon_a}{s_a}.
\end{equation*}
因此
\begin{equation*}
	E(\hat{z}_a)=z_{0a},\quad
	\operatorname{var}(\hat{z}_a)=\frac{1}{s_a^2}.
\end{equation*}
这不是定性结论，而是由正交分解直接给出的方差公式.小奇异值方向上的误差必然按 \(1/s_a^2\) 放大.

\subsection{零阶 Tikhonov 的谱闭式解}

考虑
\begin{equation*}
	Q_\lambda(\vb{\mu})
	=\frac{\alpha}{2}\left\|\tilde{\vb{y}}-\tilde{R}\vb{\mu}\right\|^2
	+\frac{\lambda}{2}\left\|\vb{\mu}\right\|^2,
	\quad \lambda>0.
\end{equation*}
这对应于
\begin{equation*}
	2G=\lambda I_M.
\end{equation*}
仍令 \(\vb{z}=Q^T\vb{\mu}\)，\(\vb{d}=U^T\tilde{\vb{y}}\).由于 \(Q\) 正交，有
\begin{equation*}
	\left\|\vb{\mu}\right\|^2
	=\left\|\vb{z}\right\|^2.
\end{equation*}
于是
\begin{equation*}
	Q_\lambda
	=\frac{\alpha}{2}\sum_{a=1}^{r}(d_a-s_az_a)^2
	+\frac{\alpha}{2}\sum_{a=r+1}^{N}d_a^2
	+\frac{\lambda}{2}\sum_{a=1}^{M}z_a^2.
\end{equation*}
对 \(z_a\) 求偏导.当 \(a\leq r\) 时，
\begin{equation*}
	\frac{\partial Q_\lambda}{\partial z_a}
	=-\alpha s_a(d_a-s_az_a)+\lambda z_a.
\end{equation*}
令其为零，得到
\begin{equation*}
	(\alpha s_a^2+\lambda)\hat{z}_a
	=\alpha s_ad_a.
\end{equation*}
因此
\begin{equation*}
	\hat{z}_a
	=\frac{\alpha s_a}{\alpha s_a^2+\lambda}d_a,
	\quad a=1,\ldots,r.
\end{equation*}
当 \(a>r\) 时，目标函数中只含 \(\lambda z_a^2/2\)，故
\begin{equation*}
	\hat{z}_a=0.
\end{equation*}
于是零阶 Tikhonov 的闭式解为
\begin{equation*}
	\hat{\vb{\mu}}_\lambda
	=\sum_{a=1}^{r}
	\frac{\alpha s_a}{\alpha s_a^2+\lambda}
	d_a\vb{q}_a.
\end{equation*}

将
\begin{equation*}
	d_a=s_az_{0a}+\epsilon_a
\end{equation*}
代入，得
\begin{equation*}
	\hat{z}_a
	=\frac{\alpha s_a^2}{\alpha s_a^2+\lambda}z_{0a}
	+\frac{\alpha s_a}{\alpha s_a^2+\lambda}\epsilon_a.
\end{equation*}
故
\begin{equation*}
	E(\hat{z}_a)-z_{0a}
	=-\frac{\lambda}{\alpha s_a^2+\lambda}z_{0a},
	\quad a\leq r,
\end{equation*}
且
\begin{equation*}
	\operatorname{var}(\hat{z}_a)
	=\left(\frac{\alpha s_a}{\alpha s_a^2+\lambda}\right)^2.
\end{equation*}
若 \(a>r\)，则
\begin{equation*}
	E(\hat{z}_a)-z_{0a}=-z_{0a},\quad
	\operatorname{var}(\hat{z}_a)=0.
\end{equation*}
这些公式完整给出了正则化带来的偏差与方差.

\subsection{一般二次惩罚的广义谱分解}

令
\begin{equation*}
	H=R^TWR.
\end{equation*}
考虑
\begin{equation*}
	Q_\lambda(\vb{\mu})
	=\frac{\alpha}{2}(\vb{y}-R\vb{\mu})^TW(\vb{y}-R\vb{\mu})
	+\frac{\lambda}{2}\vb{\mu}^TK\vb{\mu},
\end{equation*}
其中 \(K=K^T>0\).该式对应
\begin{equation*}
	2G=\lambda K.
\end{equation*}
取广义本征值问题
\begin{equation*}
	H\vb{q}_a=\rho_aK\vb{q}_a,
	\quad
	\vb{q}_a^TK\vb{q}_b=\delta_{ab}.
\end{equation*}
由于 \(K\) 正定，可以令
\begin{equation*}
	K^{1/2}\vb{q}_a=\vb{p}_a.
\end{equation*}
代入广义本征方程，得到普通对称本征值问题
\begin{equation*}
	K^{-1/2}HK^{-1/2}\vb{p}_a
	=\rho_a\vb{p}_a.
\end{equation*}
因此可以选取一组 \(K\)-正交归一的 \(\vb{q}_a\).把
\begin{equation*}
	\vb{\mu}=\sum_{a=1}^{M}z_a\vb{q}_a
\end{equation*}
代入目标函数.与 \(\vb{\mu}\) 有关的部分为
\begin{equation*}
	\frac{1}{2}\vb{\mu}^T(\alpha H+\lambda K)\vb{\mu}
	-\alpha\vb{\mu}^TR^TW\vb{y}.
\end{equation*}
利用
\begin{equation*}
	\vb{q}_a^TH\vb{q}_b
	=\rho_b\vb{q}_a^TK\vb{q}_b
	=\rho_b\delta_{ab},
\end{equation*}
得到
\begin{equation*}
	Q_\lambda
	=\frac{1}{2}\sum_{a=1}^{M}(\alpha\rho_a+\lambda)z_a^2
	-\alpha\sum_{a=1}^{M}z_a\vb{q}_a^TR^TW\vb{y}
	+\text{与 }\vb{\mu}\text{ 无关的项}.
\end{equation*}
因此
\begin{equation*}
	(\alpha\rho_a+\lambda)\hat{z}_a
	=\alpha\vb{q}_a^TR^TW\vb{y},
\end{equation*}
即
\begin{equation*}
	\hat{\vb{\mu}}
	=\sum_{a=1}^{M}
	\frac{\alpha\vb{q}_a^TR^TW\vb{y}}
	{\alpha\rho_a+\lambda}\vb{q}_a.
\end{equation*}
这就是一般正定二次惩罚下的广义谱闭式解.

\begin{notes}
	若 \(K\) 半正定而非正定，则上面的广义谱推导不能直接使用 \(K^{-1/2}\).此时仍可回到矩阵闭式解
	\begin{equation*}
		\hat{\vb{\mu}}
		=(\alpha H+\lambda K)^{-1}\alpha R^TW\vb{y},
	\end{equation*}
	但必须单独检查 \(\alpha H+\lambda K\) 是否非奇异.对于一阶差分惩罚，常数向量属于 \(K\) 的零空间，因此只要 \(H\) 在该零空间上非零，闭式解仍然唯一.
\end{notes}

\subsection{可逆、过定、欠定与秩亏的特殊情形}

若 \(N=M\)，且 \(R\) 可逆，同时不加正则化，则
\begin{equation*}
	\hat{\vb{\mu}}=R^{-1}\vb{y}.
\end{equation*}
协方差为
\begin{equation*}
	U=R^{-1}V(R^{-1})^T.
\end{equation*}
若 \(N\geq M\)，且 \(R\) 列满秩，则无正则化加权最小二乘解为
\begin{equation*}
	\hat{\vb{\mu}}
	=(R^TWR)^{-1}R^TW\vb{y}.
\end{equation*}
其协方差为
\begin{equation*}
	U=(R^TWR)^{-1}.
\end{equation*}
证明如下：
\begin{equation*}
	U
	=(R^TWR)^{-1}R^TWVWR(R^TWR)^{-1}
	=(R^TWR)^{-1}.
\end{equation*}

若 \(R\) 不列满秩，则 \(R^TWR\) 奇异.令 \(\tilde{R}=V^{-1/2}R\)，使用上一节的 SVD，所有最小二乘解具有形式
\begin{equation*}
	\vb{\mu}
	=\sum_{a=1}^{r}\frac{d_a}{s_a}\vb{q}_a
	+\sum_{a=r+1}^{M}z_a\vb{q}_a,
\end{equation*}
其中 \(z_{r+1},\ldots,z_M\) 任意.最小范数解取这些任意系数为零.

若 \(N<M\)，则即使 \(R\) 行满秩，仍有非平凡零空间
\begin{equation*}
	\ker R=\{\vb{h}:R\vb{h}=0\}.
\end{equation*}
对任意 \(\vb{h}\in\ker R\)，有
\begin{equation*}
	R(\vb{\mu}+\vb{h})=R\vb{\mu},
\end{equation*}
所以数据不能区分 \(\vb{\mu}\) 与 \(\vb{\mu}+\vb{h}\).若加入使 \(M=\alpha R^TWR+2G\) 正定的二次项，则闭式解重新唯一.

\begin{example}
	两区间分辨率模型.

	考虑
	\begin{equation*}
		R=\frac{1}{2}
		\begin{pmatrix}
			1+\varepsilon & 1-\varepsilon\\
			1-\varepsilon & 1+\varepsilon
		\end{pmatrix},
		\quad 0<\varepsilon<1.
	\end{equation*}
	取 \(V=I\).令
	\begin{equation*}
		\vb{q}_+=\frac{1}{\sqrt{2}}\binom{1}{1},
		\quad
		\vb{q}_-=\frac{1}{\sqrt{2}}\binom{1}{-1}.
	\end{equation*}
	直接计算得
	\begin{equation*}
		R\vb{q}_+=\vb{q}_+,\quad R\vb{q}_-=\varepsilon\vb{q}_-.
	\end{equation*}
	因此奇异值为
	\begin{equation*}
		s_+=1,\quad s_-=\varepsilon.
	\end{equation*}
	若 \(\vb{d}=Q^T\vb{y}\)，则无正则化解为
	\begin{equation*}
		\hat{\vb{\mu}}
		=d_+\vb{q}_++\frac{d_-}{\varepsilon}\vb{q}_-.
	\end{equation*}
	若 \(E(\vb{y})=R\vb{\mu}_0\) 且 \(\operatorname{cov}(\vb{y})=I\)，则
	\begin{equation*}
		\operatorname{var}(\hat{z}_+)=1,\quad
		\operatorname{var}(\hat{z}_-)=\frac{1}{\varepsilon^2}.
	\end{equation*}
\end{example}

\begin{example}
	一阶差分正则化矩阵.

	设
	\begin{equation*}
		S(\vb{\mu})=-\sum_{i=1}^{M-1}(\mu_i-\mu_{i+1})^2.
	\end{equation*}
	要求
	\begin{equation*}
		S(\vb{\mu})=-\vb{\mu}^TG\vb{\mu}.
	\end{equation*}
	展开得
	\begin{equation*}
		\sum_{i=1}^{M-1}(\mu_i-\mu_{i+1})^2
		=\mu_1^2+\mu_M^2
		+2\sum_{i=2}^{M-1}\mu_i^2
		-2\sum_{i=1}^{M-1}\mu_i\mu_{i+1}.
	\end{equation*}
	所以
	\begin{equation*}
		G=
		\begin{pmatrix}
			1 & -1 & 0 & \cdots & 0\\
			-1 & 2 & -1 & \ddots & \vdots\\
			0 & -1 & 2 & \ddots & 0\\
			\vdots & \ddots & \ddots & 2 & -1\\
			0 & \cdots & 0 & -1 & 1
		\end{pmatrix}.
	\end{equation*}
	对常数向量 \(\vb{e}=(1,1,\ldots,1)^T\)，有
	\begin{equation*}
		G\vb{e}=0.
	\end{equation*}
	因此一阶差分正则化不惩罚整体常数方向.若 \(R^TWR\) 在该方向上非零，则 \(M=\alpha R^TWR+2G\) 仍可正定.
\end{example}
